\uuid{Z4M5}
\exo7id{3407}
\titre{exo7 3407}
\auteur{quercia}
\organisation{exo7}
\datecreate{2010-03-10}
\isIndication{false}
\isCorrection{true}
\chapitre{Matrice}
\sousChapitre{Matrice et application linéaire}
\cpge{mpsi-mp2i.matrices-applications-lineaires.a, pcsi.matrices-applications-lineaires.a, ptsi.matrices-applications-lineaires.a, mpsi-mp2i.calcul-matriciel-systemes.e, pcsi.determinants.c, ptsi.determinants.a, bcpst1.matrices, bcpst2.applications-lineaires-matrices.c}
\module{Algèbre}
\niveau{L1}
\difficulte{}

\contenu{
\texte{
Soit $A \in \mathcal{M}_n(K)$ telle que $a_{ij} = (-1)^{n-j}C_{n-j}^{i-1}$.
}
\begin{enumerate}
    \item \question{Interpréter $A$ comme la matrice d'un endomorphisme de $ K_{n-1}[X]$.}
\reponse{$\phi(P) = (-X-1)^{n-1}P\left(-\frac 1{X+1}\right)$.}
    \item \question{En déduire $A^3$.}
\reponse{$I$.}
\end{enumerate}
}
